\section{Installationsanleitung für MySQL}

\subsection{Installation}

\subsubsection{Installation von MySQL}

Im Kursverzeichnis findest du im Ordner \myFile{32\_Datenbanken\_und\_ER-Modell} und dort
im Ordner \myFile{Installationsdateien}\footnote{oder alternativ auch hier:
\url{https://dev.mysql.com/downloads/installer/}} die Datei
\myFile{mysql-installer-web-community-8.0.46.0.msi}. Dabei handelt es sich um
ein kleines Windows-Programm welches den Download und die Installation der
eigentlichen MySQL Programme bequem erledigt.

Im Prinzip kannst du dich einfach durchklicken. Nur dann, wenn nach dem Passwort
für den root Benutzer für MySQL gefragt wird, musst du etwas eingeben: 

Bei \myPMI{Current Root Passwort} gibst du entweder nichts ein (wenn dies die
erste Installation von MySQL auf deinem Rechner ist) oder das alte
Root-Passwort. Bei \myPMI{MySQL Root Password} und \myPMI{Repeat Password} gibst
du jeweils \myUserInput{root} ein.


\subsubsection{Installation des Java-Treibers}

Um von einem Java-Programm aus auf MySQL zugreifen zu können, muss man einen speziellen Treiber
installieren, und in Eclipse einbinden (siehe einleitendes Kapitel zur
Konfiguration von Eclipse).

Im selben Ordner wie zuvor den MySQL-Installer findest auch den MySQL
JDBC-Connectors:

\myFile{mysql-connector-java-26.7.0.jar}

Dieses Java Archiv muss in euer eigenes Java-Projekt in Eclipse eingebunden
werden, damit ihr später aus eueren Java-Programmen heraus auf den MySQL-Server
zugreifen könnt. Siehe Kapitel \ref{mysql-connector-installation}.


\subsubsection{Konfiguration in Eclipse}

Die Konfiguration des DBeaver Plugins für Eclipse wird im Kapitel 1
(Eclipse als Java-Entwicklungsumgebung) beschrieben.
